\chapter{Dysarthria}
\index{Dysarthria}

\section*{Definition and Overview}
Dysarthria is a speech disorder in which the muscles used for speaking become weak, slow, uncoordinated, or difficult to control, making speech slurred, slow, or hard to understand. It is caused by damage to the brain or nerves that control the muscles of the lips, tongue, jaw, palate, and vocal cords. Dysarthria is a symptom of an underlying condition rather than a disease in itself, and its severity ranges from mild slurring to speech that is very difficult to understand. This chapter explains the causes of dysarthria, how it is diagnosed, and the therapies that help people communicate effectively.

\section*{Epidemiology}
Dysarthria is common in people with neurological conditions. It affects a large proportion of people who have had a stroke, and is also frequent in Parkinson's disease, multiple sclerosis, motor neurone disease, and cerebral palsy. The condition can develop suddenly, as after a stroke or head injury, or gradually, as in neurodegenerative disease. Dysarthria affects people of all ages and can have a significant impact on social participation and quality of life.

\section*{Aetiology and Pathophysiology}
Speech requires precise coordination of the brain, nerves, and muscles of the respiratory and vocal systems. Dysarthria results when any part of this system is damaged.
\begin{itemize}
    \item \textbf{Stroke:} damage to the motor areas of the brain disrupts the control of speech muscles
    \item \textbf{Head injury:} traumatic brain injury can damage the speech pathways
    \item \textbf{Neurodegenerative diseases:} Parkinson's disease, motor neurone disease, and multiple sclerosis progressively affect speech
    \item \textbf{Cerebral palsy:} damage before or around birth affects muscle control
    \item \textbf{Brain tumours:} tumours in the speech areas of the brain
    \item \textbf{Infections and inflammation:} conditions such as meningitis or encephalitis
    \item \textbf{Medications and toxins:} alcohol, sedatives, and some drugs can temporarily cause slurred speech
    \item \textbf{Muscle and nerve disorders:} conditions such as myasthenia gravis or muscular dystrophy
\end{itemize}

\section*{Clinical Features}
The pattern of speech difficulty depends on which part of the system is affected.
\begin{itemize}
    \item Slurred or mumbled speech that is hard to understand
    \item Slow, effortful speech with long pauses
    \item A breathy, hoarse, or strained voice
    \item Nasal-sounding speech
    \item Difficulty with volume, speaking too softly or too loudly
    \item Monotone speech with little variation in pitch
    \item Reduced clarity that worsens with fatigue
    \item Associated difficulties with chewing, swallowing, and drooling in some conditions
\end{itemize}

\section*{Differential Diagnosis}
Dysarthria should be distinguished from other communication disorders.
\begin{itemize}
    \item Aphasia, a language disorder in which the person has difficulty finding words or understanding language, despite normal speech muscles
    \item Apraxia of speech, in which the brain struggles to plan the movements for speech
    \item Voice disorders such as hoarseness or vocal cord paralysis
    \item Normal fatigue or stress-related slurring
    \item Hearing loss, which can affect speech clarity
    \item Developmental speech delays in children
\end{itemize}

\section*{When to Seek Medical Care}
Sudden onset of slurred speech requires immediate medical attention, as it may be a sign of stroke. Gradually developing dysarthria should be assessed by a GP, who can arrange investigation and referral to a neurologist and speech pathologist.

\textbf{Contact your local emergency services for:}
\begin{itemize}
    \item Sudden slurred speech, facial drooping, or weakness on one side of the body, which are signs of stroke
    \item Sudden confusion or inability to speak or understand
    \item Difficulty breathing or swallowing
    \item Collapse or loss of consciousness
\end{itemize}

\section*{Diagnosis and Investigations}
Diagnosis identifies the type of dysarthria and its underlying cause.
\begin{itemize}
    \item \textbf{History:} onset, pattern, and associated symptoms
    \item \textbf{Speech assessment:} a speech pathologist evaluates articulation, voice, resonance, and intelligibility
    \item \textbf{Neurological examination:} assessment of cranial nerves, muscle tone, and coordination
    \item \textbf{Imaging:} CT or MRI of the brain to identify stroke, tumour, or other damage
    \item \textbf{Blood tests:} to look for infection, metabolic, or inflammatory causes
    \item \textbf{Electromyography (EMG):} to assess the nerves and muscles involved in speech in some cases
\end{itemize}

\section*{Management}
Management focuses on treating the cause where possible and maximising communication.
\begin{itemize}
    \item \textbf{Speech therapy:} exercises to strengthen and coordinate the speech muscles, and strategies to improve clarity
    \item \textbf{Communication strategies:} slowing down, taking breaths, and choosing quieter environments
    \item \textbf{Communication aids:} apps, speech-generating devices, and alphabet charts for people with severe difficulty
    \item \textbf{Treating the cause:} medication for Parkinson's disease, thrombolysis for stroke, and treatment of the underlying condition
    \item \textbf{Family involvement:} educating family and friends in techniques that help them understand the person
    \item \textbf{Supporting swallowing:} assessment and management of any associated swallowing problems
\end{itemize}

\section*{Complications}
\begin{itemize}
    \item Difficulty communicating, leading to frustration and social isolation
    \item Misunderstandings and difficulty with work and daily interactions
    \item Aspiration pneumonia if swallowing is also affected
    \item Emotional distress and depression related to communication loss
    \item Reduced participation in social and family life
\end{itemize}

\section*{Prognosis and Outcomes}
The outlook for dysarthria depends on the cause. Dysarthria after a stroke or head injury often improves substantially with speech therapy as the brain recovers. In progressive conditions such as motor neurone disease, speech typically declines, and communication aids become increasingly important. With appropriate therapy and support, most people find ways to communicate effectively and maintain their connections with others.

\section*{Prevention and Self-Care}
\begin{itemize}
    \item Recognise the signs of stroke and call for help immediately, as early treatment limits damage
    \item Attend speech therapy regularly and practise the recommended exercises
    \item Speak slowly and take time to communicate
    \item Reduce background noise when talking
    \item Tell listeners what helps you communicate, such as needing them to face you
    \item Seek referral to a speech pathologist early for assessment and advice
\end{itemize}

\section*{Sources and Further Reading}
The following reputable sources provide further information for healthcare students and patients seeking reliable, current advice about dysarthria.
\begin{itemize}
    \item Speech Pathology Australia. \textit{Speech and language information.} https://www.speechpathologyaustralia.org.au/
    \item NHS. \textit{Dysarthria: causes, symptoms, and treatment.} https://www.nhs.uk/conditions/dysarthria/
    \item Mayo Clinic. \textit{Speech disorders: causes and treatment.} https://www.mayoclinic.org/diseases-conditions/dysarthria/
    \item MedlinePlus. \textit{Dysarthria.} https://medlineplus.gov/ency/article/007470.htm
    \item American Speech-Language-Hearing Association. \textit{Dysarthria information.} https://www.asha.org/
    \item National Institute on Deafness and Other Communication Disorders. \textit{Speech disorders.} https://www.nidcd.nih.gov/
\end{itemize}
